% =========================================================
% Subsets of R — Notes
% Chapter: functions
% Section: subsets-real-line
% Volume III · Analysis
% =========================================================

\section{Subsets of $\mathbb{R}$}
\label{sec:subsets-of-R}

% ---------------------------------------------------------
% Toolkit
% ---------------------------------------------------------

\begin{tcolorbox}[title={Toolkit --- Subsets of $\mathbb{R}$}]
\begin{itemize}
  \item Cluster point (= limit point = accumulation point): every
    $\varepsilon$-neighbourhood contains a \emph{distinct} point of $X$.
    Adherent point: every $\varepsilon$-neighbourhood contains \emph{some}
    point of $X$ (witness may equal $x$).
  \item Adherent = cluster $\cup$ isolated. Closure $\overline{X}$ =
    set of all adherent points.
  \item Closed $\iff$ $\overline{X} = X$ $\iff$ sequentially closed
    (limits of sequences in $X$ stay in $X$).
  \item Heine--Borel: $X$ closed and bounded $\iff$ every sequence in
    $X$ has a subsequence converging to a limit in $X$.
  \item Every point of an interval is a cluster point of that interval.
    This justifies intervals as the natural domains for function limits.
\end{itemize}
\end{tcolorbox}

\begin{remark*}[Terminological note]
The terms \emph{cluster point}, \emph{limit point}, and
\emph{accumulation point} are synonyms. Bartle and Lebl use
\emph{cluster point}; Rudin uses \emph{limit point}; Tao uses both.
The term \emph{adherent point} is Tao's notation for the weaker notion
that includes isolated points. These notes follow the convention:
cluster point = strict notion (distinct witness required); adherent
point = weak notion (any witness, including $x$ itself).
\end{remark*}

% ---------------------------------------------------------
% Definition — epsilon-Neighbourhood
% ---------------------------------------------------------

\begin{tcolorbox}[title={Definition --- $\varepsilon$-Neighbourhood}]
\begin{definition}[$\varepsilon$-Neighbourhood]
\label{def:epsilon-neighbourhood}
Let $x \in \mathbb{R}$ and $\varepsilon > 0$. The
\textbf{$\varepsilon$-neighbourhood} of $x$ is
\[
  V_\varepsilon(x) \;:=\; (x - \varepsilon,\; x + \varepsilon)
  \;=\; \{\, y \in \mathbb{R} : |y - x| < \varepsilon \,\}.
\]
The \textbf{deleted $\varepsilon$-neighbourhood} of $x$ is
\[
  V_\varepsilon^*(x) \;:=\; V_\varepsilon(x) \setminus \{x\}
  \;=\; \{\, y \in \mathbb{R} : 0 < |y - x| < \varepsilon \,\}.
\]
\end{definition}
\end{tcolorbox}

\begin{remark*}[Standard quantified statement]
\[
  y \in V_\varepsilon(x) \;\iff\; |y - x| < \varepsilon.
  \qquad
  y \in V_\varepsilon^*(x) \;\iff\; 0 < |y - x| < \varepsilon.
\]
\end{remark*}

\begin{remark*}[Interpretation]
The deleted neighbourhood $V_\varepsilon^*(x)$ is the operative set
in the definition of a function limit: the condition $0 < |x - c| <
\delta$ means exactly $x \in V_\delta^*(c)$. The $0 <$ excludes the
base point $c$ itself, which is the defining structural feature
separating limits from continuity.
\end{remark*}

% ---------------------------------------------------------
% Definition — Cluster Point
% ---------------------------------------------------------

\begin{tcolorbox}[title={Definition --- Cluster Point}]
\begin{definition}[Cluster Point]
\label{def:cluster-point-R}
Let $X \subseteq \mathbb{R}$ and let $x \in \mathbb{R}$. The point
$x$ is a \textbf{cluster point} of $X$ (also called a \textbf{limit
point} or \textbf{accumulation point} of $X$) if every
$\varepsilon$-neighbourhood of $x$ contains at least one point of
$X$ distinct from $x$:
\[
  \forall\,\varepsilon > 0,\quad
  \exists\, y \in X \setminus \{x\},\quad |y - x| < \varepsilon.
\]
\end{definition}
\end{tcolorbox}

\begin{remark*}[Standard quantified statement]
\[
  x \text{ is a cluster point of } X
  \;\iff\;
  \forall\,\varepsilon > 0,\;
  \exists\, y \in X \setminus \{x\} :\;
  |y - x| < \varepsilon
  \;\iff\;
  x \text{ is adherent to } X \setminus \{x\}.
\]
\end{remark*}

\begin{remark*}[Negated quantified statement]
\[
  x \text{ is not a cluster point of } X
  \;\iff\;
  \exists\,\varepsilon > 0,\;
  \forall\, y \in X \setminus \{x\} :\;
  |y - x| \geq \varepsilon.
\]
\end{remark*}

\begin{remark*}[Interpretation]
A cluster point is a point that $X$ approaches from arbitrarily close
by distinct witnesses. A cluster point need not belong to $X$: $0$ is
a cluster point of $(0,1)$ but $0 \notin (0,1)$. A finite set and
$\mathbb{N}$ have no cluster points at all. The concept is exactly
what is needed to make the function limit non-vacuous: $\lim_{x \to c}
f(x)$ requires $c$ to be a cluster point of $\operatorname{dom}(f)$
so that the condition $0 < |x-c| < \delta$ can actually be satisfied.
\end{remark*}

% ---------------------------------------------------------
% Theorem — Sequential Characterization of Cluster Points
% ---------------------------------------------------------

\begin{theorem}[Sequential Characterization of Cluster Points]
\label{thm:cluster-point-sequential}
Let $A \subseteq \mathbb{R}$ and let $c \in \mathbb{R}$. Then $c$
is a cluster point of $A$ if and only if there exists a sequence
$(a_n) \subseteq A \setminus \{c\}$ with $\lim_{n \to \infty} a_n
= c$.

\smallskip
\noindent
\hyperref[prf:cluster-point-sequential]{\textit{Go to proof.}}
\end{theorem}

\begin{remark*}[Standard quantified statement]
\[
  c \text{ cluster point of } A
  \;\iff\;
  \exists (a_n)_{n=1}^\infty \subseteq A \setminus \{c\} :\;
  a_n \to c.
\]
\end{remark*}

\begin{remark*}[Interpretation]
This bridges the $\varepsilon$-$\delta$ definition with the sequential
language used throughout the chapter. The forward direction uses the
$\varepsilon = 1/n$ construction to build the witness sequence. The
reverse direction is immediate since the sequence witnesses the
$\varepsilon$-$\delta$ condition directly.
\end{remark*}

% ---------------------------------------------------------
% Definition — Adherent Point and Isolated Point
% ---------------------------------------------------------

\begin{tcolorbox}[title={Definition --- Adherent Point and Isolated Point}]
\begin{definition}[Adherent Point; Isolated Point]
\label{def:adherent-point-R}
Let $X \subseteq \mathbb{R}$ and let $x \in \mathbb{R}$.

The point $x$ is an \textbf{adherent point} of $X$ if every
$\varepsilon$-neighbourhood of $x$ contains at least one point of
$X$:
\[
  \forall\,\varepsilon > 0,\quad
  \exists\, y \in X :\quad |y - x| \leq \varepsilon.
\]

A point $x \in X$ is an \textbf{isolated point} of $X$ if $x \in X$
and $x$ is not a cluster point of $X$; equivalently,
\[
  \exists\,\varepsilon > 0 :\; V_\varepsilon(x) \cap X = \{x\}.
\]
\end{definition}
\end{tcolorbox}

\begin{remark*}[Standard quantified statement]
\[
  x \text{ adherent to } X
  \;\iff\;
  \forall\,\varepsilon > 0,\;
  V_\varepsilon(x) \cap X \neq \emptyset.
\]
\end{remark*}

\begin{remark*}[Cluster vs.\ adherent]
The sole difference is whether the witness $y$ may equal $x$. Every
cluster point is adherent; an isolated point is adherent but not a
cluster point. The partition is:
\[
  \text{Adherent points of } X
  = \text{Cluster points of } X
  \;\cup\;
  \text{Isolated points of } X.
\]
\end{remark*}

% ---------------------------------------------------------
% Definition — Interior, Exterior, Boundary
% ---------------------------------------------------------

\begin{definition}[Interior, Exterior, Boundary]
\label{def:interior-exterior-boundary-R}
Let $X \subseteq \mathbb{R}$ and $x \in \mathbb{R}$.

$x$ is an \textbf{interior point} of $X$ if some neighbourhood of
$x$ lies entirely in $X$:
\[
  \exists\,\varepsilon > 0 :\; V_\varepsilon(x) \subseteq X.
\]
The \textbf{interior} of $X$ is $X^\circ := \{x : x \text{ interior
to } X\}$.

$x$ is a \textbf{boundary point} of $X$ if every neighbourhood of
$x$ meets both $X$ and $X^c$:
\[
  \forall\,\varepsilon > 0 :\;
  V_\varepsilon(x) \cap X \neq \emptyset
  \;\wedge\;
  V_\varepsilon(x) \cap X^c \neq \emptyset.
\]
The \textbf{boundary} of $X$ is $\partial X$.
\end{definition}

\begin{remark*}[Standard quantified statement]
\[
\begin{aligned}
  x \in X^\circ
  &\;\iff\; \exists\,\varepsilon > 0,\;
    \forall\, y :\; |y-x| < \varepsilon \Rightarrow y \in X. \\
  x \in \partial X
  &\;\iff\; \forall\,\varepsilon > 0,\;
    V_\varepsilon(x) \cap X \neq \emptyset
    \;\wedge\;
    V_\varepsilon(x) \cap X^c \neq \emptyset.
\end{aligned}
\]
\end{remark*}

\begin{remark*}[Interpretation]
The three sets $X^\circ$, $\partial X$, and $\mathrm{Ext}(X) :=
(X^c)^\circ$ partition $\mathbb{R}$. $X$ is open iff $X = X^\circ$
(contains no boundary points). $X$ is closed iff $X \supseteq
\partial X$ (contains all its boundary points). The closure
$\overline{X} = X^\circ \cup \partial X$.
\end{remark*}

% ---------------------------------------------------------
% Definition — Closure
% ---------------------------------------------------------

\begin{tcolorbox}[title={Definition --- Closure}]
\begin{definition}[Closure of a Subset of $\mathbb{R}$]
\label{def:closure-R}
Let $X \subseteq \mathbb{R}$. The \textbf{closure} of $X$ is the set
of all adherent points of $X$:
\[
  \overline{X}
  \;:=\;
  \bigl\{\, x \in \mathbb{R} \;\bigm|\;
    \forall\varepsilon > 0,\;
    \exists\, y \in X :\;
    |x - y| \leq \varepsilon
  \,\bigr\}.
\]
\end{definition}
\end{tcolorbox}

\begin{remark*}[Standard quantified statement]
\[
  x \in \overline{X}
  \;\iff\;
  \forall\varepsilon > 0,\;
  V_\varepsilon(x) \cap X \neq \emptyset.
\]
\end{remark*}

\begin{remark*}[Canonical examples]
$\overline{(a,b)} = [a,b]$;\quad
$\overline{\mathbb{Q}} = \mathbb{R}$;\quad
$\overline{\mathbb{Z}} = \mathbb{Z}$;\quad
$\overline{\{1/n : n \in \mathbb{N}\}} = \{1/n : n \in \mathbb{N}\}
\cup \{0\}$.
\end{remark*}

% ---------------------------------------------------------
% Lemma — Elementary Properties of Closures
% ---------------------------------------------------------

\begin{lemma}[Elementary Properties of Closures]
\label{lem:closure-elementary}
Let $X, Y \subseteq \mathbb{R}$. Then:
\begin{enumerate}[label=(\roman*)]
  \item $X \subseteq \overline{X}$.
  \item $\overline{X \cup Y} = \overline{X} \cup \overline{Y}$.
  \item $\overline{X \cap Y} \subseteq \overline{X} \cap \overline{Y}$.
  \item $X \subseteq Y \Rightarrow \overline{X} \subseteq \overline{Y}$.
\end{enumerate}

\smallskip
\noindent
\hyperref[prf:closure-elementary]{\textit{Go to proof.}}
\end{lemma}

\begin{remark*}[Standard quantified statement]
\[
  \forall X, Y \subseteq \mathbb{R}:\quad
  X \subseteq \overline{X};\quad
  \overline{X \cup Y} = \overline{X} \cup \overline{Y};\quad
  \overline{X \cap Y} \subseteq \overline{X} \cap \overline{Y};\quad
  X \subseteq Y \Rightarrow \overline{X} \subseteq \overline{Y}.
\]
\end{remark*}

\begin{remark*}[Failure of equality in (iii)]
$\overline{X \cap Y} = \overline{X} \cap \overline{Y}$ fails in general.
Take $X = (0,1)$, $Y = (1,2)$: then $\overline{X \cap Y} =
\overline{\emptyset} = \emptyset$ but $\overline{X} \cap \overline{Y}
= [0,1] \cap [1,2] = \{1\}$.
\end{remark*}

% ---------------------------------------------------------
% Definition — Closed Set
% ---------------------------------------------------------

\begin{tcolorbox}[title={Definition --- Closed Subset of $\mathbb{R}$}]
\begin{definition}[Closed Subset of $\mathbb{R}$]
\label{def:closed-set-R}
A set $E \subseteq \mathbb{R}$ is \textbf{closed} if $E$ contains
all its adherent points:
\[
  E \text{ is closed}
  \;\iff\;
  \overline{E} = E.
\]
\end{definition}
\end{tcolorbox}

\begin{remark*}[Standard quantified statement]
\[
  E \text{ closed}
  \;\iff\;
  \forall x \in \mathbb{R} :\;
  \bigl(\forall\varepsilon > 0,\;
    \exists\, y \in E :\; |x-y| \leq \varepsilon\bigr)
  \;\Rightarrow\; x \in E.
\]
\end{remark*}

\begin{remark*}[Negated quantified statement]
\[
  E \text{ not closed}
  \;\iff\;
  \exists x \notin E :\;
  \forall\varepsilon > 0,\;
  V_\varepsilon(x) \cap E \neq \emptyset.
\]
\end{remark*}

\begin{remark*}[Interpretation]
Closedness is the condition that the set contains its own boundary.
Canonical examples: $[a,b]$ is closed; $(a,b)$ is not (its boundary
points $a,b$ lie outside it); $\mathbb{R}$ and $\emptyset$ are both
closed; $\mathbb{Q}$ is not closed since $\overline{\mathbb{Q}} =
\mathbb{R} \neq \mathbb{Q}$.
\end{remark*}

% ---------------------------------------------------------
% Corollary — Sequential Characterization of Closed Sets
% ---------------------------------------------------------

\begin{corollary}[Closed iff Sequentially Closed]
\label{cor:closed-iff-seq-limits}
Let $X \subseteq \mathbb{R}$. Then $X$ is closed if and only if
every convergent sequence $(a_n) \subseteq X$ has its limit in $X$:
\[
  X \text{ closed}
  \;\iff\;
  \forall (a_n) \subseteq X,\;
  a_n \to x \;\Rightarrow\; x \in X.
\]

\smallskip
\noindent
\hyperref[prf:closed-iff-seq-limits]{\textit{Go to proof.}}
\end{corollary}

\begin{remark*}[Interpretation]
A set is closed precisely when it cannot be escaped by convergent
sequences. The forward direction: if $x$ is adherent to $X$ then
some sequence in $X$ converges to $x$; closedness pulls $x$ into $X$.
The reverse direction: every adherent point has a sequence witnessing
adherence; the sequential hypothesis forces it into $X$.
\end{remark*}

% ---------------------------------------------------------
% Corollary — Every Point of an Interval Is a Cluster Point
% ---------------------------------------------------------

\begin{corollary}[Every Point of an Interval Is a Cluster Point]
\label{cor:interval-all-limit-points}
Let $I$ be an interval. Then every element of $I$ is a cluster point
of $I$.

\smallskip
\noindent
\hyperref[prf:interval-all-limit-points]{\textit{Go to proof.}}
\end{corollary}

\begin{remark*}[Interpretation]
This justifies intervals as the natural domains for function limits.
The condition $\lim_{x \to c} f(x)$ requires $c$ to be a cluster
point of the domain; since every point of an interval satisfies this,
no additional hypotheses are needed when working on intervals.
\end{remark*}

% ---------------------------------------------------------
% Definition — Bounded Subset
% ---------------------------------------------------------

\begin{definition}[Bounded Subset of $\mathbb{R}$]
\label{def:bounded-set-R}
A subset $X \subseteq \mathbb{R}$ is \textbf{bounded} if
$X \subseteq [-M, M]$ for some $M > 0$:
\[
  X \text{ bounded}
  \;\iff\;
  \exists M > 0,\;
  \forall x \in X :\; |x| \leq M.
\]
\end{definition}

\begin{remark*}[Standard quantified statement]
\[
  X \text{ unbounded}
  \;\iff\;
  \forall M > 0,\;
  \exists x \in X :\; |x| > M.
\]
\end{remark*}

\begin{remark*}[Interpretation]
A bounded set fits inside some finite symmetric interval. For a
nonempty bounded set, both $\sup X$ and $\inf X$ exist in $\mathbb{R}$
by the Axiom of Completeness.
\end{remark*}

% ---------------------------------------------------------
% Theorem — Heine-Borel
% ---------------------------------------------------------

\begin{theorem}[Heine--Borel Theorem for $\mathbb{R}$]
\label{thm:heine-borel}
Let $X \subseteq \mathbb{R}$. The following are equivalent:
\begin{enumerate}[label=(\alph*)]
  \item $X$ is closed and bounded.
  \item Every sequence $(a_n) \subseteq X$ has a subsequence
    converging to a limit $L \in X$.
\end{enumerate}

\smallskip
\noindent
\hyperref[prf:heine-borel]{\textit{Go to proof.}}
\end{theorem}

\begin{remark*}[Standard quantified statement]
\[
  \bigl(X \text{ closed} \wedge X \text{ bounded}\bigr)
  \;\iff\;
  \forall (a_n) \subseteq X,\;
  \exists (a_{n_k}),\;
  \exists L \in X :\;
  a_{n_k} \to L.
\]
\end{remark*}

\begin{remark*}[Interpretation]
Closed and bounded together characterize the subsets of $\mathbb{R}$
on which every sequence must eventually settle. Boundedness prevents
sequences from escaping to infinity; closedness prevents limits from
escaping through boundary gaps. The forward direction: bounded gives
Bolzano--Weierstrass; closed keeps the limit inside $X$. The reverse
direction: if $X$ were unbounded, the sequence $x_n$ with $|x_n| > n$
would have no convergent subsequence; if $X$ were not closed, an
adherent point outside $X$ would be the limit of a sequence in $X$
with no limit in $X$.
\end{remark*}

% ---------------------------------------------------------
% Definition — True Near a Point
% ---------------------------------------------------------

\begin{definition}[True Near a Point]
\label{def:true-near}
Let $x_0 \in \mathbb{R}$, let $f$ be a function defined near $x_0$
(except possibly at $x_0$ itself). A property $P(f(x))$ is
\textbf{true near $x_0$} if there exists $\delta > 0$ such that
$P(f(x))$ holds for all $x$ with $0 < |x - x_0| < \delta$.
\end{definition}

\begin{remark*}[Standard quantified statement]
\[
  P(f(x)) \text{ true near } x_0
  \;\iff\;
  \exists\,\delta > 0,\;
  \forall x :\; 0 < |x - x_0| < \delta \Rightarrow P(f(x)).
\]
\end{remark*}

\begin{remark*}[Interpretation]
The limit definition $\lim_{x \to x_0} f(x) = L$ can be restated:
for every $\varepsilon > 0$, the property $|f(x) - L| < \varepsilon$
is true near $x_0$. This compact notation makes the logical structure
of limit arguments transparent.
\end{remark*}
